\section{Interpolation by single-source distance functions}
\label{int:section}

All distances in this section are intrinsic to one connected component
$Y$ of a closed band $X_J$.  We write $d_J$ for this distance.  The width
of $J$ is denoted by $a$, and its midpoint by $H_*$.  The active column
graph and its radial coordinate $b_i=d_J(o,(i,H_*))$ are as in
Section~\ref{geo:section}; we choose $o$ at height $H_*$.  The choices of
$o$ and of all portals below are made before any random parameters in
this section are sampled.

The construction will repeatedly use signed point-distance formulas as
local data. We therefore approximate a function of Lipschitz constant
strictly below one by formulas $c\pm d_J(k,\cdot)$, while controlling where the
selected formula changes. Each formula has a fixed source and sign, a
random additive constant, and a name. When the construction reuses a
formula, it retains that name. Stability on a small ball means that one
name is selected throughout the ball, not merely that the values agree
at the interfaces.

\begin{theorem}[Thin-band interpolation]\label{int:interpolation}
Let $0<\zeta<1$, let $E>0$, and set
\begin{equation}\label{int:constants}
 \begin{gathered}
 \delta=E/8,\qquad a_0=10^{-6}\zeta\delta,\qquad
 q=\zeta\delta/20,\\
 \rho_0=\zeta\delta/1000,\qquad
 M=\lceil15/\zeta\rceil+3,\qquad N=6(M+1).
 \end{gathered}
\end{equation}
Suppose $0<a\le a_0$, and let $p:Y\to\R$ be
$(1-\zeta)$-Lipschitz for $d_J$.  There is a random function
$Q:Y\to\R$, obtained using finitely many independent scalar noises
uniform on $[q,2q]$, with the following properties.
\begin{enumerate}
\item For every realization, $Q$ respects all switches, is
$1$-Lipschitz for $d_J$, and satisfies $|Q(v)-p(v)|\le6\delta<E$.
\item Its pieces use finitely many named atoms of the form
\begin{equation}\label{int:atom-form}
 A(v)=C+\varepsilon d_J(k,v),\qquad \varepsilon\in\{-1,1\},\quad k\in Y.
\end{equation}
Each name specifies a source $k$ and sign $\varepsilon$ fixed before the
noises are sampled; the intercept $C$ is random. Whenever the construction
reuses an atom, it retains its name.
\item For a fixed $x\in Y$ and $0<\rho\le\rho_0$, there is a list
$\mathcal A(x,\rho)$ of at most $N$ atom names and sources, determined
without the noises, that contains the name of every atom selected by $Q$ on
$B_{d_J}(x,\rho)$.  Except on an event of probability at most
\begin{equation}\label{int:probability}
 \frac{N^2}{q}(8\rho+\eta),
\end{equation}
one named atom is selected throughout this ball.
\end{enumerate}
Here $\eta=0$ for continuous uniform noises.  Alternatively, each noise
may be uniform on an equally spaced grid including both endpoints of
$[q,2q]$, and $\eta$ in \eqref{int:probability} may be the largest grid
spacing.  The same conclusions hold for any set of points reached
from $x$ by band paths of length at most $\rho$, including the points
and representatives encountered on those paths.
\end{theorem}

\subsection{The tree of superlevel components}

For every integer $n\ge0$, take the connected components of the
subgraph on columns with $b_i\ge n\delta$.  Its nonempty components
are called nodes of depth $n$.  The sole node of depth zero is the
root.  A node of depth $n+1$ is contained in exactly one node of depth
$n$, its parent.  Because there are finitely many columns and all
$b_i$ are finite, there are only finitely many nodes.  The layer
belonging to a node $\nu$ of depth $n$ consists of the full columns in
that node satisfying
\[
 n\delta\le b_i<(n+1)\delta.
\]
A node can have an empty layer; it is still retained in the tree.
Every column has exactly one layer node.

The layers specify where a node's formula will be used.  The formula itself
will be defined on all of $Y$.  Near an entrance to a child node, the child
and parent formulas will agree exactly; this will allow us to paste their
restrictions across layer boundaries without changing values or names.

For a positive-depth node, choose its set $K_\nu$ of one or two
entrance portals using Lemma~\ref{geo:portals}.  The portals are
midheight points on entrance columns.  At the root set $K_\nu=\{o\}$.
The radial and portal estimates of Lemma~\ref{geo:portal-levels} give
the following facts, which also apply to the root with this convention:
\begin{align}
 \min_{k\in K_\nu}d_J(k,v)&\le\delta+220a
 &&\text{if $v$ is in the layer of $\nu$ or is a child portal},
 \label{int:near-portal}\\
 d_J(k,v)&\ge\delta-10a
 &&\text{if $k\in K_\nu$ and $v$ is a child portal}.
 \label{int:separated-portals}
\end{align}
For a positive-depth node of depth $n$, every $k\in K_\nu$ also
satisfies
\begin{equation}\label{int:portal-radial}
 n\delta\le d_J(o,k)\le n\delta+2a.
\end{equation}
Every point on an entrance column is within $200a$ of $K_\nu$.
These are statements about $d_J$, so paths witnessing them may leave
the corresponding node.

\subsection{The formulas and their deterministic estimates}

At each node we first approximate the input by positive distance cones
based at its portals. We then constrain this tentative formula so that
it agrees with the parent near every entrance. The strict Lipschitz
bound on the input will limit how many generations an accurate atom
can be inherited through.

For every portal $k\in K_\nu$, sample an independent noise
$\alpha_{\nu k}\in[q,2q]$. At every nonroot node also sample
$\alpha'_{\nu k},\alpha''_{\nu k}\in[q,2q]$. All these noises are
mutually independent. Define formulas on the whole of $Y$, recursively
down the tree. The tentative formula is
\[
 T_\nu(v)=\min_{k\in K_\nu}
       \bigl(p(k)+d_J(k,v)+\alpha_{\nu k}\bigr).
\]
At the root put $F_\nu=T_\nu$. This is $1$-Lipschitz, because
distance functions and their finite minima are $1$-Lipschitz.

Now let $\nu$ have parent $\pi$, whose formula is already defined
and $1$-Lipschitz. To preserve the parent at the interfaces, we will
constrain $T_\nu(v)$ to an interval containing $F_\pi(v)$ that
collapses to this single value near $K_\nu$. Define two envelopes
of cones based at the parent values:
\begin{align*}
 A_\nu(v)&=\min_{k\in K_\nu}
     \bigl(F_\pi(k)+d_J(k,v)-\alpha'_{\nu k}\bigr),\\
 B_\nu(v)&=\max_{k\in K_\nu}
     \bigl(F_\pi(k)-d_J(k,v)+\alpha''_{\nu k}\bigr).
\end{align*}
If $d_J(k,v)<q/2$ for a portal $k$, the term with source $k$
in the first envelope is below $F_\pi(v)$, and its term in the
second is above $F_\pi(v)$. Indeed the parent changes by at most
$d_J(k,v)$, whereas each noise is at least $q$.
Thus $A_\nu(v)<F_\pi(v)<B_\nu(v)$ near every portal, even if
the other portal's cones have different values. The desired interval is
\[
 [L_\nu(v),U_\nu(v)]
 =\bigl[\min(F_\pi(v),B_\nu(v)),
        \max(F_\pi(v),A_\nu(v))\bigr].
\]
It always contains the parent value and reduces to that value in
these portal neighborhoods. Clamp the tentative formula to this interval:
\[
 F_\nu(v)=\min\bigl(U_\nu(v),\max(L_\nu(v),T_\nu(v))\bigr).
\]
All the formulas remain $1$-Lipschitz on $Y$, since finite minima and
maxima preserve this property. This completes the recursive construction.

The new atoms of a nonroot node are
\begin{align}
 p(k)+d_J(k,v)+\alpha_{\nu k},\qquad
 F_\pi(k)+d_J(k,v)-\alpha'_{\nu k},\qquad
 F_\pi(k)-d_J(k,v)+\alpha''_{\nu k}.
 \label{int:three-types}
\end{align}
The root has its single tentative atom. An intercept $F_\pi(k)$ in
\eqref{int:three-types} is a scalar, not an instruction to inherit the
atom selected at $k$. Each expression in \eqref{int:three-types} is a
new atom with source $k$. In contrast, selecting the branch
$F_\pi(v)$ inherits its atom and its name.

Fix a total priority order on names in which all smaller-depth names
precede larger-depth names. In every minimum or maximum, resolve
equal values in favor of the earliest name. Multiple appearances of
one inherited atom have the same name. This completely specifies a
named atom selected by every formula at every point.

\begin{lemma}[Accuracy and equality collars]\label{int:deterministic}
Every node formula $F_\nu$ has error at most $5\delta$ at every portal
of a child of $\nu$, and error at most $6\delta$ on its own layer.
At a nonroot node with parent $\pi$, if
\[
 \dist_{d_J}(v,K_\nu)\le c_0:=\zeta\delta/50,
\]
then $F_\nu(v)=F_\pi(v)$, with the same selected atom name.
\end{lemma}

\begin{proof}
Since $c_0<q/2$, the portal-neighborhood calculation above gives
$U_\nu(v)=L_\nu(v)=F_\pi(v)$ in the stated collar. The priority
convention retains the parent atom even if $T_\nu(v)=F_\pi(v)$.

For accuracy, the Lipschitz bound of $p$ gives, at every point,
\begin{equation}\label{int:tentative-bounds}
 p(v)+\zeta\dist_{d_J}(v,K_\nu)+q
 \le T_\nu(v)
 \le p(v)+(2-\zeta)\dist_{d_J}(v,K_\nu)+2q.
\end{equation}
In the situations of \eqref{int:near-portal}, its upper bound is
less than $p(v)+3\delta$, by \eqref{int:constants}.  This proves both
accuracy assertions at the root.

Proceed by induction.  At a nonroot node, the inductive hypothesis
gives $|F_\pi(k)-p(k)|\le5\delta$ for all its own portals.  Therefore,
writing $r_\nu(v)=\dist_{d_J}(v,K_\nu)$,
\begin{align}
 A_\nu(v)&\ge p(v)-5\delta+\zeta r_\nu(v)-2q,
 \label{int:upper-atom-bound}\\
 B_\nu(v)&\le p(v)+5\delta-\zeta r_\nu(v)+2q.
 \label{int:lower-atom-bound}
\end{align}
At a child portal, \eqref{int:separated-portals} and
$a\le10^{-6}\zeta\delta$ imply
$\zeta r_\nu(v)\ge2q$.  Hence
$U_\nu(v)\ge p(v)-5\delta$ and
$L_\nu(v)\le p(v)+5\delta$.  Clamping a number in
$[p(v),p(v)+3\delta]$ to an interval with these two properties gives
a value in $[p(v)-5\delta,p(v)+5\delta]$.

On the own layer, the same argument without the favorable distance
term gives $U_\nu(v)\ge p(v)-5\delta-2q$ and
$L_\nu(v)\le p(v)+5\delta+2q$.  Since $2q\le\delta/10$, the
clamped value has error at most $5.1\delta<6\delta$.
\end{proof}

Define $Q$ on a representative $(i,H)$ by the formula of the unique
layer node of column $i$.  Two columns joined by a switch have radial
values differing by at most $2a<\delta$, so their layer depths are
equal or consecutive.  At equal depths their nodes coincide.  At
consecutive depths the upper node is a child of the lower node, and
the upper column is an entrance column.  Its switch point is within
$200a<c_0$ of an upper portal, so Lemma~\ref{int:deterministic} gives
the same value and atom name from both representatives.  Applying this
to every switch also gives consistency under transitive identifications.
Thus $Q$ is defined on $Y$.  It is $1$-Lipschitz on each column and
respects switches; summing along band paths shows that it is
$1$-Lipschitz for $d_J$.  Its accuracy follows from
Lemma~\ref{int:deterministic}.

\subsection{One deterministic ancestor chain for a small ball}

The pasted function is now well defined and has the required approximation
error.  To control changes of its selected formula, we next show that a
small ball sees one node formula and only boundedly many of its ancestors.
This step uses the geometry and the strict Lipschitz bound of the input;
the random comparison estimate comes afterward.

\begin{lemma}\label{int:common-formula}
For $x\in Y$ and $0<\rho\le\rho_0$, one node $\nu$, determined
without the interpolation noises, satisfies $Q=F_\nu$ on
$B_{d_J}(x,\rho)$, including equality of selected atom names.
If $n$ is the depth of $\nu$, every point of this ball has a
representative on a column with $b_i\ge n\delta$.
\end{lemma}

\begin{proof}
The closed ball is connected by geodesics through its center.  Include
all representatives of all its points in its column projection.  For
any two such representatives their radial values differ by at most
$2\rho+2a<\delta$, using the radial comparison with $d_J(o,\cdot)$.
Thus their layer depths take either a single value or two consecutive
values $n,n+1$.  The projected graph is connected, and all its columns
have $b_i\ge n\delta$.  They therefore belong to one superlevel node
$\nu$ at depth $n$.  In the single-depth case this proves the result.

Otherwise choose a represented point of layer depth $n$ in the ball.
For any represented point $v$ of depth $n+1$, concatenate its geodesic
to the center with a geodesic to the chosen point.  This is a band
path of length at most $2\rho$.  Its last crossing from depth $n$
to depth $n+1$, when traversed toward $v$, enters the child node
containing the representative of $v$.  The entry column is an entrance
column of that child.  Hence $v$ is within
$2\rho+200a<c_0$ of a portal of that child.  The equality collar of
Lemma~\ref{int:deterministic} identifies its formula and its selected
atom with those of $\nu$ at $v$.  This proves the assertion for every
representative.  All the choices in this argument depend only on the
band geometry and the ball.
\end{proof}

\begin{lemma}[Bounded depth of a selected atom]\label{int:depth}
For the node $\nu$ supplied by Lemma~\ref{int:common-formula}, every
atom selected by $F_\nu$ on the ball originates in $\nu$ or one of
its last $M$ ancestors.  Consequently at most $N$ atom names can be
selected there.
\end{lemma}

\begin{proof}
We first prove an obstruction to passing a parent atom through a
distant node $\mu$.  Suppose
\begin{equation}\label{int:far-portals}
 \zeta d_J(k,v)>13\delta\quad\text{for every }k\in K_\mu.
\end{equation}
Equations \eqref{int:tentative-bounds}--\eqref{int:lower-atom-bound}
give
\[
 T_\mu(v)>p(v)+13\delta,\qquad
 A_\mu(v)>p(v)+7\delta,\qquad
 B_\mu(v)<p(v)-7\delta.
\]
If the parent value belonged to $[p(v)-6\delta,p(v)+6\delta]$, then
$L_\mu(v)=B_\mu(v)$ and $U_\mu(v)=A_\mu(v)$, so
\[
 F_\mu(v)=\min(A_\mu(v),T_\mu(v))>p(v)+7\delta.
\]
In particular, a value within $6\delta$ of $p(v)$ cannot be obtained
at this node by selecting its parent atom.  Notice that the lower
bound on the tentative formula is required here as well as the two
clamping bounds.

Let $n$ be the depth of $\nu$.  If $n\le M$, every ancestor is already
in the asserted list.  Otherwise consider its ancestor $\mu$ of
depth $n-M\ge1$.  For a point $v$ of the ball,
Lemma~\ref{int:common-formula}, radial comparison, and
\eqref{int:portal-radial} give
\[
 d_J(k,v)\ge M\delta-3a\qquad(k\in K_\mu).
\]
The definition of $M$ implies \eqref{int:far-portals}.  An atom
originating strictly earlier than $\mu$ could be selected by $F_\nu(v)$
only by being inherited through every intervening node.  All these
inherited values would equal $F_\nu(v)=Q(v)$, whose error is at most
$6\delta$.  In particular $\mu$ would select its parent atom with
such a value, contradicting the preceding paragraph.  This excludes
all earlier origins.  Each of the remaining $M+1$ nodes contributes
at most six atoms, proving the count $N$.
\end{proof}

\subsection{Conditional comparisons and probability bounds}

We have a list of at most $N$ possible names on the ball, fixed before the
noises are sampled.  Their intercepts can depend on earlier noises.
The next lemma identifies an independent noise in each comparison, so
this dependence does not require a union bound over the whole node tree.

\begin{lemma}[Fresh-noise comparison]\label{int:near-tie}
For any two distinct named atoms in the deterministic list of
Lemma~\ref{int:depth}, any fixed $x\in Y$, and $t\ge0$,
\begin{equation}\label{int:near-tie-bound}
 \Pr\bigl(|A(x)-A'(x)|\le t\bigr)\le(2t+\eta)/q.
\end{equation}
The same bound holds when one atom is replaced by any real quantity
independent of all current interpolation noises.
\end{lemma}

\begin{proof}
For two atoms originating at different depths, use the later atom's
own noise.  Its intercept, whether $p(k)$ or $F_\pi(k)$, does not
depend on that noise, because parent formulas use only earlier nodes.
The earlier atom does not depend on it either.  Thus, after conditioning
on every other noise, the difference has the form $c+\varepsilon\alpha$
with $\varepsilon\in\{-1,1\}$ and $\alpha$ uniform on $[q,2q]$.
For two atoms at the same node their own noises are different and
their parent intercepts depend on neither, so the same argument
applies.  The candidate list lies on one ancestor chain, so these
are all cases.  An independent comparison quantity can likewise be
included in the conditioning.

For a continuous uniform noise an interval of length $2t$ has
probability at most $2t/q$.  For a grid with spacing $\eta'\le\eta$,
it contains at most $2t/\eta'+1$ grid points; division by the total
number $q/\eta'+1$ gives at most $(2t+\eta')/q$.
\end{proof}

\begin{proof}[Completion of Theorem~\ref{int:interpolation}]
Take the deterministic list furnished by Lemma~\ref{int:depth}.
The union bound and Lemma~\ref{int:near-tie}, with $t=4\rho$, show
that the probability of a pair satisfying
$|A(x)-A'(x)|\le4\rho$ is at most
$N^2(8\rho+\eta)/q$.  On the complementary event, two atoms in the
list cannot be equal anywhere on the ball, since both are
$1$-Lipschitz.  The continuous function $Q=F_\nu$ selects a value
from this finite list at every point.  Where the list values are
distinct its selected name is locally constant: continuity and the
positive minimum gap to the other finitely many values give a
neighborhood with the same selection.  A ball is connected by paths
through its center, so the name is constant throughout the ball.
This proves \eqref{int:probability}.  The path version follows by
applying the ball assertion to the paths, all of whose points lie
in the ball.
\end{proof}

\begin{corollary}[Comparison with other branches]\label{int:comparisons}
Condition on arbitrary prior data. Let $W$ be the union of a nonempty
family of band paths issuing from a fixed point $x$, each of length at
most $\rho$. The set $W$ and these paths are fixed by the prior data.
Suppose that $Q_1,\ldots,Q_s$ are interpolation constructions, each
on a band component containing all these paths. Their geometries,
inputs, and portals are fixed by the prior data, and their fresh noise
families are mutually independent. The bands and their intrinsic metrics
may differ; each of the specified paths must be permitted in every band.
Denote the constants of $Q_i$ in \eqref{int:constants} by
$q_i,N_i,\rho_{0,i}$.

Allow inputs $c_i+\varepsilon_iQ_i$, with
$\varepsilon_i\in\{-1,1\}$ and $c_i$ fixed by the prior data.
Also allow one fixed branch $f$ and $t$ branches $g_j+\beta_j$,
where $f$ and the $g_j$ are fixed functions whose restrictions to each
specified path, parametrized by arclength, are $L$-Lipschitz. The
$\beta_j$ are fresh independent uniform noises on intervals of length
$w_j>0$, independent also of the interpolation noises. All functions respect the switches used by the paths.
Set
\[
 L_* =\max(1,L),\qquad
 S=1+t+\sum_{i=1}^sN_i,\qquad
 q_* =\min\bigl(\{q_i:1\le i\le s\}\cup
                         \{w_j:1\le j\le t\}\bigr).
\]
Assume that there is at least one randomized input and that
$\rho\le\min_i\rho_{0,i}$ when $s>0$. Any fixed expression formed
from these inputs by finitely many minima and maxima selects one
branch and, for a selected interpolation, one atom throughout $W$,
except with conditional probability at most
\begin{equation}\label{int:comparison-probability}
 \frac{S^2}{q_*}(8L_*\rho+\eta).
\end{equation}
Continuous noises have $\eta=0$; for equally spaced endpoint grids
$\eta$ can be the largest spacing of any noise used here. The
expression uses each interpolation with a single chosen sign and
shift; repeated appearances of that same input are given the same
name.
\end{corollary}

\begin{proof}
Every specified path lies in $B_{d_{J_i}}(x,\rho)$ for each interpolation
band $J_i$. Use the deterministic bounded atom list for $Q_i$ on that
ball, with its chosen sign and shift. Adjoin the shifted $g_j$ branches
and, if present, $f$.
For pairs from one interpolation use Lemma~\ref{int:near-tie}.
For pairs from different interpolations condition on all noises
except the own noise of an atom in one interpolation; independence
of the constructions makes the other atom independent of that noise.
The same argument applies to a shifted $g_j$ or the fixed branch
$f$. There is at most one branch without a fresh noise. Therefore
every comparison of distinct names at $x$ has the interval bound
$(2u+\eta)/q_*$ at tolerance $u$.

Take $u=4L_*\rho$ and sum over at most $S^2$ pairs. Off the event
in \eqref{int:comparison-probability}, no two listed values agree
anywhere on $W$: their changes along each path from $x$ are bounded
by $L_*\rho$. The expression is a continuous selection from the
resulting finite list along each path. Its selected name is therefore
constant on each path and, because all paths begin at $x$, throughout
$W$.
\end{proof}

\begin{remark}\label{int:conditional-scope}
The conditional independence requirements in
Corollary~\ref{int:comparisons} are substantive.  In particular, a
branch defined using current interpolation noises cannot be treated
as a fixed comparison branch, and two arbitrary fixed branches need
not have a small probability of meeting.  The later applications use
one old branch, plus independently shifted height or ordered branches,
after conditioning on the bands, old charts, groups, and admissions.
All choices of geometric portals and interpolation inputs precede
the noises being tested.  With this conditioning, the displayed
bounds are uniform in all prior choices and in the size of the
column graph.
\end{remark}
